% !TeX program = pdflatex
\documentclass[11pt,a4paper]{article}

% The original course sheets are pdfLaTeX documents.  Keep their classical
% Computer Modern / CM-Super typography, but inline the shared preamble so
% that this file also compiles outside the site repository.
\usepackage{cmap}
\usepackage[T2A]{fontenc}
\usepackage[utf8]{inputenc}
\usepackage[english,bulgarian]{babel}
\usepackage{fontawesome5}
\usepackage{amsmath,amssymb,mathtools}
\usepackage[a4paper,top=20mm,bottom=21mm,left=21mm,right=21mm,headheight=16pt]{geometry}
\usepackage{booktabs,longtable,tabularx,array,multicol}
\usepackage[table]{xcolor}
\usepackage{enumitem}
\usepackage[ruled,vlined]{algorithm2e}
\usepackage{fancyhdr}
\usepackage{hyperref}
\usepackage[expansion=false]{microtype}
\usepackage{titlesec}
\usepackage{pifont}
\usepackage{wasysym}
\usepackage{xparse}

% Keep T2A text scalable even on installations where CM-Super is present but
% its maps have not yet been registered by updmap.
%\ifdefined\pdfmapfile
%  \pdfmapfile{+cm-super-minimal-t2a.map}
%  \pdfmapfile{+cm-super-extended-t2a.map}
%\fi

\definecolor{siteblue}{HTML}{173B8F}
\definecolor{sitegold}{HTML}{F2C94C}
\definecolor{muted}{HTML}{64748B}
\definecolor{rulegray}{HTML}{B8C0CC}
\definecolor{softgray}{HTML}{F4F5F7}

\hypersetup{
  unicode=true,
  pdfencoding=auto,
  colorlinks=true,
  linkcolor=siteblue,
  urlcolor=siteblue,
  citecolor=siteblue,
  pdfauthor={Martin Minchev},
  pdftitle={Вероятности и статистика: седмица 1 -- Комбинаторика, 2026/2027}
}

\pagestyle{fancy}
\fancyhf{}
\lhead{2026/2027}
\rhead{\nouppercase{\leftmark}}
\cfoot{\thepage}
\renewcommand{\headrulewidth}{0.4pt}
\renewcommand{\sectionmark}[1]{\markboth{#1}{#1}}

\titleformat{\section}{\normalfont\large\bfseries}{\thesection.}{0.55em}{}
\titleformat{\subsection}{\normalfont\normalsize\bfseries}{}{0pt}{}
\titlespacing*{\section}{0pt}{1.2em}{0.55em}
\titlespacing*{\subsection}{0pt}{0.85em}{0.3em}
\newcommand{\sectionbreak}{\clearpage}

\setlist[itemize]{leftmargin=1.55em,itemsep=0.18em,topsep=0.25em}
\setlist[enumerate]{leftmargin=2.2em,itemsep=0.35em,topsep=0.3em}
\newenvironment{enumeratea}{\begin{enumerate}[label=(\alph*),itemsep=0.28em,topsep=0.25em]}{\end{enumerate}}
\setlength{\parindent}{0pt}
\setlength{\parskip}{0.32em}
\setlength{\emergencystretch}{3em}
\renewcommand{\arraystretch}{1.15}
\raggedbottom

% Compatibility names for the icon language of the existing exercise bank.
% \newcommand{\faBook}{\ding{110}}
% \newcommand{\faStar}{\ding{72}}
% \NewDocumentCommand{\faSmile}{O{}}{\smiley}
% \NewDocumentCommand{\faPencil}{s}{\ding{45}}

\newcommand{\E}{\mathbb E}
\renewcommand{\P}{\mathbb P}
\newcommand{\Pbb}{\mathbb P}
\newcommand{\Ebb}{\mathbb E}
\newcommand{\Var}{\operatorname{Var}}
\newcommand{\Cov}{\operatorname{Cov}}
\newcommand{\Cor}{\operatorname{Cor}}
\newcommand{\Ber}{\operatorname{Ber}}
\newcommand{\Bin}{\operatorname{Bin}}
\newcommand{\Poi}{\operatorname{Poi}}
\newcommand{\Geom}{\operatorname{Geom}}
\newcommand{\Exp}{\operatorname{Exp}}
\newcommand{\Unif}{\operatorname{U}}
\newcommand{\ind}{\mathbf 1}
\newcommand{\iid}{\mathrel{\mathrm{iid}}}

\newcounter{problem}[section]
\renewcommand{\theproblem}{\thesection.\arabic{problem}}
\newcommand{\problemstart}[2]{%
  \refstepcounter{problem}%
  \par\medskip\noindent
  #1\if\relax\detokenize{#1}\relax\else\,\fi
  \textbf{Задача \theproblem.}%
  \if\relax\detokenize{#2}\relax\else\ \emph{(#2)}\fi
  \enspace\ignorespaces
}
\NewDocumentEnvironment{problem}{O{}}{\problemstart{}{#1}}{\par}
\NewDocumentEnvironment{problembook}{O{}}{\problemstart{\faBook}{#1}}{\par}
\NewDocumentEnvironment{problemstar}{O{}}{\problemstart{\faStar}{#1}}{\par}
\NewDocumentEnvironment{probleminteresting}{O{}}{\problemstart{\faSmile[regular]}{#1}}{\par}
\NewDocumentEnvironment{problempencil}{O{}}{\problemstart{\faPencil*}{#1}}{\par}
\NewDocumentEnvironment{problemstarpencil}{O{}}{\problemstart{\faStar\,\faPencil*}{#1}}{\par}
\NewDocumentEnvironment{researchtask}{O{}}{\problemstart{\faRobot}{#1}}{\par}

\newcommand{\aitag}[1]{{\footnotesize\sffamily\color{muted}AI--#1}}
\newcommand{\projectcheck}[1]{%
  \par\smallskip\noindent\colorbox{softgray}{%
    \parbox{\dimexpr\linewidth-2\fboxsep}{\small\textbf{Проектен checkpoint.} #1}}\par\smallskip
}
\newcommand{\leancheck}[1]{%
  \par\smallskip\begingroup\sloppy\noindent{\small\textbf{Lean micro-task.} #1}\par\endgroup\smallskip
}

% Before publishing, fill in the actual model names and versions used.
% "Latest versions" alone does not identify a model/version under SU Art. 27.
\newcommand{\chatgptmodel}{\textit{[модел и версия]}}
\newcommand{\claudemodel}{\textit{[модел и версия]}}

\begin{document}

\setcounter{tocdepth}{1}
\tableofcontents

\subsection*{Бележки}

Всяка тема съдържа \emph{основни задачи}, които покриват минимума,
необходим за съответната тема; \emph{допълнителни задачи} за разнообразяване
и надграждане на основните; и \emph{предизвикателства}, които са свързани
с темата, но излизат извън изискванията на стандартните уводни курсове
по вероятности и статистика.

Появяват се и следните означения:
\begin{itemize}
  \item[\faBook] задача, представяща основни обекти от курса;
  \item[\faStar] предизвикателна задача;
  \item[\faSmile] интересен пример;
  \item[\faPencil*] задача от контролно във ФМИ;
  \item[\faRobot] задача за самостоятелно изследване. Разрешени са всички
  помощни средства, включително литература, софтуер и AI, като се очаква да проверите резултатите, както и
  да ги разберете и можете да обясните.
\end{itemize}


{\small
\textbf{Бележка за използването на AI.}
За форматиране, редактиране и уточняване на отделни формулировки
са използвани последните публично достъпни версии на ChatGPT и Claude (в момента на писане съответно sol 6.1  и opus 5.5).
Отговорността за окончателния вариант и математическата му коректност
е на съставителя.
}

\section{Комбинаторика}

\subsection*{Основни задачи}

\begin{problembook}
За $n,k\ge 1$ припомнете си:
\begin{enumeratea}
  \item колко решения има уравнението $x_1+\cdots+x_n=k$, когато
  $x_1,\dots,x_n$ са положителни цели числа? А когато са неотрицателни цели числа?
  \item формулирайте принципа за включване и изключване.
\end{enumeratea}
\end{problembook}

\begin{problembook}
Нека $M$ е множество с $n$ елемента. По колко начина можем да изберем:
\begin{enumeratea}
  \item $k$ различни елемента от $M$;
  \item наредена $k$-орка от различни елементи на $M$;
  \item наредена $k$-орка от елементи на $M$, като повторенията са позволени?
\end{enumeratea}
\end{problembook}

\begin{problem}
По колко начина можем да разпределим всички $k$ различими частици в $n$
различни клетки, ако:
\begin{enumeratea}
  \item всяка клетка съдържа най-много една частица;
  \item клетките съдържат произволен брой частици;
  \item няма празна клетка?
\end{enumeratea}
Отговорете на същите въпроси, когато частиците са неразличими. Посочете
и в кои случаи отговорът е нула.
\end{problem}

\begin{problem}
Пет различими топки се разпределят в три различни кутии $A,B,C$. Намерете броя
на разпределенията, при които:
\begin{enumeratea}
  \item $A$ е празна;
  \item само $A$ е празна;
  \item точно една кутия е празна;
  \item поне една кутия е празна;
  \item няма празна кутия.
\end{enumeratea}
\end{problem}

\begin{problemstar}
Разгледайте следните въпроси.
\begin{enumeratea}
  \item По колко начина можем да стигнем от $(0,0)$ до $(n,n)$ в квадратната
  решетка с единични стъпки надясно и нагоре? А ако пътят не преминава над диагонала?
  \item Колко са правилните скобови изрази с $n$ отварящи и $n$ затварящи
  скоби? Постройте явна биекция между тях и пътищата от втората част на
  предходната точка.
  \item Намерете прости асимптотични еквиваленти на броя на всички пътища
  и на броя на пътищата, които не преминават над диагонала, при $n\to\infty$.
\end{enumeratea}
\end{problemstar}

\subsection*{Допълнителни задачи}

\begin{problem}
Колко четирицифрени числа могат да се напишат с цифрите $1,2,3,4,5$, ако:
\begin{enumeratea}
  \item не се допуска повторение;
  \item повторение се допуска;
  \item повторение не се допуска и числото е нечетно?
\end{enumeratea}
\end{problem}

\begin{problem}
По колко начина може да се избере четиричленна делегация от 12 кандидати,
сред които $A,B,C,D$ са различни, ако:
\begin{enumeratea}
  \item няма ограничения;
  \item $A$ и $B$ не трябва да участват заедно;
  \item $C$ и $D$ или участват и двамата, или не участва нито един?
\end{enumeratea}
\end{problem}

\begin{problem}
Докажете комбинаторно, а след това и алгебрично, че
\[
\sum_{k=0}^n\binom nk=2^n
\qquad
\text{и}
\qquad
\binom{2n}{n}=\sum_{k=0}^n\binom nk^2.
\]
\end{problem}

\begin{problem}
Нека $p(n)$ е броят на разбиванията на $n$ на различни положителни цели числа,
а $q(n)$ е броят на разбиванията му на нечетни положителни цели числа, като
повторенията са позволени и редът на събираемите няма значение. Например за $n=6$:
\[
p(6)=4:\ 6,\ 5+1,\ 4+2,\ 3+2+1; \qquad
q(6)=4:\ 5+1,\ 3+3,\ 3+1+1+1,\ 1+1+1+1+1+1.
\]
\begin{center}
\begin{tabular}{c|cccccccc}
  $n$    & 1 & 2 & 3 & 4 & 5 & 6 & 7 & 8 \\ \hline
  $p(n)$ & 1 & 1 & 2 & 2 & 3 & 4 & 5 & 6 \\
  $q(n)$ & 1 & 1 & 2 & 2 & 3 & 4 & 5 & 6
\end{tabular}
\end{center}
Като сравните пораждащите функции, докажете, че $p(n)=q(n)$ за всяко $n$.
\end{problem}

\subsection*{Предизвикателства и самостоятелно изследване}

\begin{problemstar}[Derangements]
Нека $D_n$ е броят на пермутациите $\sigma\in S_n$, за които
$\sigma(i)\ne i$ за всяко $i$.
\begin{enumeratea}
  \item Докажете, че
  \[
  D_n=n!\sum_{j=0}^{n}\frac{(-1)^j}{j!}.
  \]
  \item Докажете комбинаторно, че за $n\ge 2$
  \[
  D_n=(n-1)(D_{n-1}+D_{n-2}),
  \]
  където $D_0=1$ и $D_1=0$.
  \item Докажете, че за $n\ge 1$ числото $D_n$ е най-близкото цяло число до
  $n!/e$, без да използвате числено приближение на $e$.
\end{enumeratea}
\end{problemstar}

\begin{problemstar}[Огърлици и лемата на Бърнсайд]
Нека $X=\{0,1\}^n$ е множеството на двоичните думи с дължина $n$, а
$G=\{r^0,r^1,\dots,r^{n-1}\}\cong\mathbb Z_n$ е групата на цикличните
завъртания, $r(x_1x_2\dots x_n)=x_nx_1\dots x_{n-1}$. \emph{Огърлица}
наричаме орбита на $G$ в $X$, т.е. отъждествяваме две думи, ако едната се
получава от другата чрез циклично завъртане. Например
$0011,\,0110,\,1100,\,1001$ са една огърлица, а $0101,\,1010$ са друга.
Завъртанията не променят броя на единиците, така че $G$ действа и върху
\[
X_k=\{x\in X:\ x\text{ има точно }k\text{ единици}\}.
\]
С $N(n,k)$ означаваме броя на орбитите в $X_k$, т.е. броя на огърлиците с
$n$ мъниста и точно $k$ единици.

За $g\in G$ нека
\[
\operatorname{Fix}(g)=\{x\in X_k:\ g(x)=x\},
\]
а $X_k/G$ е множеството на орбитите. Използвайте лемата на Бърнсайд
(Cauchy--Frobenius)
\[
\#(X_k/G)=\frac1{|G|}\sum_{g\in G}|\operatorname{Fix}(g)|.
\]
Обяснете защо завъртане $g$, чиито цикли върху позициите
$\{1,\dots,n\}$ имат дължина $d$, може да фиксира дума с $k$ единици
само ако $d\mid k$. Докажете, че в този случай
\[
|\operatorname{Fix}(g)|=\binom{n/d}{k/d},
\]
а ако $d\nmid k$, то $\operatorname{Fix}(g)=\emptyset$. Докажете също,
че за всяко $d\mid n$ има точно $\varphi(d)$ завъртания, чиито цикли имат
дължина $d$. Изведете, че за $1\le k\le n-1$
\[
N(n,k)=\frac1n\sum_{d\mid\gcd(n,k)}\varphi(d)\binom{n/d}{k/d},
\]
където $\varphi$ е функцията на Ойлер. Като проверка намерете $N(12,4)$.
% Отговор: N(12,4) = (495 + 15 + 6)/12 = 43.
\end{problemstar}

\begin{researchtask}[Случайни функции и циклични върхове]
Нека $f:[n]\to[n]$ е избрана равномерно измежду всички $n^n$ функции, $G_f$ е
насоченият граф с ребра $i\to f(i)$, а $K_n$ е броят на върховете, лежащи на
цикъл в $G_f$. Можете да използвате AI за предложени решения, а вие да сте
техен \emph{рецензент}: оценява се проверката, а не генерирането.
\begin{enumeratea}
  \item Намерете чрез софтуер разпределението $\Pbb(K_n=k)$ за $n\le 7$.
  Кодът може да е генериран от AI, но го изпълнете и проверете:
  $\sum_k\#\{K_n=k\}=n^n$, а случая $n=3$ проверете и ръчно.
  \item Намерете формула за $\#\{f:K_n=k\}$ с доказателство.
  По желание поискайте предложение от AI.
  (i) Сравнете я с таблицата от т. 1. (ii) Запишете доказателството със свои
  думи, разделяйки избора на цикличните върхове, пермутацията върху тях и
  насочената гора с корени в тях. (iii) Дайте собствено доказателство на
  формулата $k\,n^{n-k-1}$ за броя на горите с $k$ зададени корена.
  Ако използвате AI, проверете дали той действително я е доказал, и дайте
  доказателство с метод, различен от неговия.
  (iv) Намерете и проверете числено броя на функциите с $K_n=k$ и точно $c$
  цикъла.
  \item Изведете $\Ebb K_n$ по два начина: от разпределението и чрез линейност,
  като пресметнете $\Pbb(\text{върхът }1\text{ е цикличен})$. Докажете
  тъждеството между двата израза. С точната формула пресметнете $\Ebb K_n$ за
  $n=10,10^2,10^3,10^4,10^5$ и изследвайте числено дали
  \[
  \frac{\Ebb K_n}{\sqrt n}\longrightarrow\sqrt{\frac{\pi}{2}}.
  \]
  Обяснете откъде идва константата $\sqrt{\pi/2}$. По желание изследвайте и
  по-финото асимптотично твърдение
  \[
  \Ebb K_n-\sqrt{\frac{\pi n}{2}}\longrightarrow-\frac13.
  \]
  Нека $T_n$ е броят на различните върхове в траекторията
  $1,f(1),f^2(1),\dots$ до първото повторение. Съвпадат ли $K_n$ и $T_n$ по
  разпределение, или имат само еднакви математически очаквания?
\end{enumeratea}
% \textbf{Проверка.} Представете таблица с колони \emph{твърдение / метод на
% проверка (собствено изчисление, собствено доказателство, източник) / резултат}
% и поне пет реда. Включете поне едно грешно или непълно твърдение на AI или
% списък на опитите да го провокирате: разпределението за $n=8$ без код;
% константата $-\tfrac13$; съвпадат ли разпределенията на $K_n$ и $T_n$.
% Ако сте използвали AI, приложете разговора. При устната защита трябва да можете да докажете
% формулата $k\,n^{n-k-1}$.
% За n=3: 9, 12, 6.
% Броят в т. 2(iv) е C(n,k)c(k,c)k n^{n-k-1}, където c(k,c) е беззнаково
% число на Стирлинг от първи род.
% E K_n = sum_{k=1}^n (n)_k/n^k = sqrt(pi n/2) - 1/3 + O(n^{-1/2}).
% P(K_n=k) = P(T_n=k) = ((n)_k/n^k) - ((n)_{k+1}/n^{k+1}).
\end{researchtask}

\noindent\textbf{Lean.} Нека $A,B\subseteq\Omega$, а
индикаторите приемат стойности в $\mathbb N$. Дефинирайте индикатора чрез
\[
\ind_A(x)=
\begin{cases}
1,&x\in A,\\
0,&x\notin A,
\end{cases}
\]
и докажете в Lean 4 с Mathlib, че за всяко $x\in\Omega$
\[
\ind_{A\cup B}(x)+\ind_{A\cap B}(x)=\ind_A(x)+\ind_B(x).
\]


\section{Условна вероятност, независимост и Бейс}

\subsection*{Основни задачи}

\begin{problembook}
Хвърлят се независимо два стандартни зара.
\begin{enumeratea}
  \item Каква е вероятността сумата да е по-малка от 8, ако се знае, че е
  нечетна? Независими ли са събитията ,,сумата е по-малка от 8`` и
  ,,сумата е нечетна``?
  \item Независими ли са събитията ,,сумата е 7`` и ,,първият зар показва 3``?
  А събитията ,,сумата е 8`` и ,,първият зар показва 3``?
\end{enumeratea}
% (a) 2/3; не: P(AB)=1/3, P(A)P(B)=7/24.
% (b) да: 1/36=(1/6)(1/6); не: 1/36 != (5/36)(1/6).
\end{problembook}
\begin{problem}
Нека $A$ и $B$ са събития с $\Pbb(A)>0$ и $0<\Pbb(B)<1$.
\begin{enumeratea}
  \item Докажете, че ако $A$ и $B$ са независими, то независими са и
  $A$ и $\overline{B}$, а също $\overline{A}$ и $\overline{B}$.
  \item Докажете, че следните четири условия са еквивалентни:
  \[
  \Pbb(A\cap B)>\Pbb(A)\Pbb(B),\quad
  \Pbb(A\mid B)>\Pbb(A),\quad
  \Pbb(B\mid A)>\Pbb(B),\quad
  \Pbb(A\mid B)>\Pbb(A\mid \overline{B}).
  \]
  \item Ако $A$ е независимо от $B$ и $A$ е независимо от $C$, следва ли,
  че $A$ е независимо от $B\cap C$? А от $B\cup C$?
\end{enumeratea}
% (a) P(A \cap B^c) = P(A) - P(A \cap B) = P(A)(1 - P(B)).
% (b) Първите три са деления на P(AB) > P(A)P(B) с P(B), съотв. P(A);
%     за четвъртото: P(A) = P(A|B)P(B) + P(A|B^c)P(B^c) е претеглено средно
%     на P(A|B) и P(A|B^c) с положителни тегла.
% (c) Не и в двата случая: A -- първата монета е ези, B -- втората е ези,
%     C -- двете са еднакви (вж. задачата за независимост по двойки):
%     P(A ∩ B ∩ C) = 1/4 != (1/2)(1/4), P(A ∩ (B ∪ C)) = 1/4 != (1/2)(3/4).
\end{problem}

\begin{problembook}
От тесте с 52 карти се теглят последователно три карти без връщане.
Намерете вероятността:
\begin{enumeratea}
  \item всичките три карти да са аса;
  \item третата карта да е асо;
  \item втората карта да е асо, ако първата не е.
\end{enumeratea}
% (a) (4/52)(3/51)(2/50)=1/5525; (b) 1/13; (c) 4/51.
\end{problembook}

\begin{problembook}
В кутия има три монети: две честни и една с ези от двете страни.
Избираме равномерно монета и я хвърляме. При фиксирана монета
последователните хвърляния са независими.
\begin{enumeratea}
  \item Пада се ези. Каква е вероятността монетата да е фалшивата?
  \item Хвърляме същата монета отново и пак се пада ези. Намерете новата
  вероятност по два начина: като обновите отговора от предходната точка,
  приет за нова априорна вероятност, и директно от двете хвърляния.
  \item Каква е вероятността третото хвърляне да е ези, ако първите две
  са били ези?
\end{enumeratea}
% (a) 1/2; (b) 2/3 по двата начина; (c) 5/6.
\end{problembook}

\begin{problempencil}[К1, СЕМ 2023]
В урна има 4 жълти, 5 зелени и 2 сини камъчета. Избираме равномерно камъче.
Ако то е жълто, хвърляме един стандартен зар; ако е зелено -- два;
ако е синьо -- три. Хвърлянията са независими.
Ако сумата е 3, каква е вероятността камъчето да е било зелено?
% (5/18)/(4/6+5/18+2/216)=30/103.
\end{problempencil}

\begin{problembook}
При болен човек тестът дава положителен резултат с вероятност $99\%$
(чувствителност), а при здрав човек дава отрицателен резултат с вероятност
$99\%$ (специфичност). Заболяването се среща при $0.5\%$ от населението.
Избираме равномерно човек и го тестваме.
\begin{enumeratea}
  \item Каква е вероятността човек с положителен резултат действително
  да е болен?
  \item При какво разпространение на заболяването тази вероятност е
  точно $1/2$?
  \item Човекът се тества повторно и резултатът отново е положителен.
  Ако двата теста са условно независими при дадено здравно състояние
  и имат същите чувствителност и специфичност, каква е вероятността
  да е болен? Кога това допускане за независимост е съмнително?
\end{enumeratea}
% (a) 99/298 ≈ 0.3322; (b) 1%; (c) 9801/10000=0.9801.
\end{problembook}


\subsection*{Допълнителни задачи}

\begin{probleminteresting}[Kahneman, Thinking, Fast and Slow]
В град със синя и зелена таксиметрова компания през нощта става катастрофа.
$85\%$ от такситата са зелени, а $15\%$ са сини. Приемете, че вероятността
едно такси да участва в катастрофа не зависи от компанията.

Свидетел твърди, че таксито е било синьо. Последвала експертиза показва,
че той разпознава правилно цвета с вероятност $80\%$, както при синьо,
така и при зелено такси. Каква е вероятността таксито наистина да е било синьо?
% 0.12/(0.12+0.17)=12/29 ≈ 0.4138.
\end{probleminteresting}

\begin{problem}
Три жетона имат страни съответно бяло--бяло, черно--черно и бяло--черно.
Избираме равномерно жетон и го поставяме със случайно избрана страна нагоре,
като двете му страни са еднакво вероятни.
Ако видим бяла горна страна, каква е вероятността скритата страна също
да е бяла?
% 2/3.
\end{problem}

\begin{probleminteresting}[Monty Hall]
Зад една от три врати, избрана равномерно, има автомобил.
Избирате врата, без да знаете къде е автомобилът. Водещият знае къде е той
и винаги отваря друга празна врата, след което ви предлага да смените
избора си с останалата затворена врата.
\begin{enumeratea}
  \item Трябва ли да смените избора? Зависи ли вероятността да спечелите при
  стратегия ,,винаги сменям`` от начина, по който водещият избира коя
  врата да отвори, когато има две възможности?
  \item Избрали сте врата 1, а водещият е отворил врата 3. Когато има избор
  между вратите 2 и 3, той отваря врата 3 с вероятност $p\in[0,1]$.
  Намерете условната вероятност автомобилът да е зад врата 2.
  За кои $p$ си струва да смените?
\end{enumeratea}
% (a) При смяна вероятността е 2/3. Равномерният избор на водещия между
% две празни врати не е необходим за тази безусловна вероятност.
% (b) 1/(1+p) ∈ [1/2, 1]; при p = 1/2 е 2/3. Смяната никога не е по-лоша.
% Проверка: (1+p)/3 * 1/(1+p) + (2-p)/3 * 1/(2-p) = 2/3.
\end{probleminteresting}

\begin{probleminteresting}[Birthday paradox]
Приемете, че рождените дни на различните хора са независими и равномерно
разпределени между 365 дни. Какъв е най-малкият брой хора, за който
вероятността поне двама да имат един и същ рожден ден е по-голяма от $1/2$?
Използвайте правилото за умножение. Обсъдете допусканията в модела.
% 23: prod_{i=0}^{22}(1-i/365) ≈ 0.4927.
\end{probleminteresting}

\begin{problem}
Хвърлят се независимо две честни монети. Нека $A$ е събитието
,,първата монета показва ези``, $B$ -- ,,втората монета показва ези``,
а $C$ -- ,,двете монети показват едно и също``.
Докажете, че $A,B,C$ са независими по двойки, но не и в съвкупност.
Независими ли са $A$ и $C$ при условие $B$?
% P(ABC)=1/4 != 1/8.
% При условие B: P(AC|B)=1/2, P(A|B)P(C|B)=1/4.
\end{problem}

\begin{probleminteresting}[Парадокс на Симпсън]
Следната таблица съдържа реални данни за два метода при лечение на бъбречни
камъни (Charig et al., \emph{BMJ} 292, 1986; Julious \& Mullee, \emph{BMJ} 309, 1994):
\begin{center}
\small
\begin{tabular}{@{}lrr@{}}
\toprule
 & Лечение A & Лечение B \\
\midrule
Малки камъни & $93\%\ (81/87)$ & $87\%\ (234/270)$ \\
Големи камъни & $73\%\ (192/263)$ & $69\%\ (55/80)$ \\
Общо & $78\%\ (273/350)$ & $83\%\ (289/350)$ \\
\bottomrule
\end{tabular}
\end{center}
Обяснете защо сравнението между двата метода се обръща при обединяване
на групите. Може ли само от тези данни да заключим кое лечение е по-добро?
% A има по-висок наблюдаван процент на успех във всяка група.
% При A обаче делът на пациентите с големи камъни е много по-висок.
% Таблицата сама по себе си не установява причинно превъзходство.
\end{probleminteresting}

\begin{probleminteresting}[Reservoir sampling]
Нека $A$ е масив с $n\ge3$ различни елемента, индексирани от 1.
Изпълнява се следният алгоритъм, като случайните избори в различните
итерации са независими:

\begin{algorithm}[H]
\caption{Reservoir sampling за един елемент}
$sample \gets A[1]$\;
\For{$i \gets 2$ \KwTo $n$}{
  $j \gets$ случайно цяло число, равномерно разпределено
  в $\{1,\dots,i\}$\;
  \If{$j = 1$}{
    $sample \gets A[i]$\;
  }
}
\Return{$sample$}
\end{algorithm}

Каква е вероятността накрая $sample$ да е равно на $A[1]$? На $A[3]$?
На $A[n]$? Докажете, че алгоритъмът избира равномерно елемент от масива.
% За всеки елемент вероятността е 1/n.
% За A[3]: (1/3) prod_{i=4}^n (1-1/i)=1/n.
\end{probleminteresting}

\begin{problem}
Разполагаме с $n\ge1$ урни, всяка с по $m$ бели и $k$ черни топки,
където $m+k>0$. От първата прехвърляме равномерно избрана топка във втората;
след добавянето ѝ избираме равномерно топка от текущото съдържание на втората
и я прехвърляме в третата, и т.н.
Намерете вероятността топка, изтеглена равномерно накрая от последната урна,
да е бяла.
% m/(m+k) за всяко n.
% p_1=m/(m+k), p_i=(m+p_{i-1})/(m+k+1).
\end{problem}


\subsection*{Предизвикателства и самостоятелно изследване}

\begin{problemstar}[Скрита обща причина]
Нека $E$ е събитие с $\Pbb(E)=q$, където $0<q<1$, а $A$ и $B$ са
условно независими при дадена стойност на $\ind_E$. Нека
\[
\Pbb(A\mid \ind_E=e)=a_e
\quad\text{и}\quad
\Pbb(B\mid \ind_E=e)=b_e,
\qquad
\text{за }e\in\{0,1\}.
\]
\begin{enumeratea}
  \item Докажете, че
  \[
  \Cov(\ind_A,\ind_B)
  =q(1-q)(a_1-a_0)(b_1-b_0).
  \]
  Кога тази ковариация е нула? Обяснете защо за индикаторите това е
  еквивалентно на независимостта на $A$ и $B$.

  \item При $\Pbb(A\cap B)>0$ намерете $\Pbb(E\mid A\cap B)$.

  \item Нека
  \[
  q=0.10,\qquad
  a_0=0.02,\quad a_1=0.20,\qquad
  b_0=0.03,\quad b_1=0.25.
  \]
  Пресметнете $\Pbb(A)$, $\Pbb(B)$ и $\Cov(\ind_A,\ind_B)$.
  Намерете и сравнете
  \[
  \Pbb(E),\qquad
  \Pbb(E\mid A),\qquad
  \Pbb(E\mid B),\qquad
  \Pbb(E\mid A\cap B).
  \]
  Проверете, че
  \[
  \Pbb(A\cap B)=0.00554
  \qquad\text{и}\qquad
  \Pbb(E\mid A\cap B)=\frac{250}{277}.
  \]
  Дайте интерпретация на резултатите, като разглеждате $E$ като скрита
  обща причина, а $A$ и $B$ като две наблюдавани улики за нея.
\end{enumeratea}
% P(A)=0.038, P(B)=0.052, Cov(ind_A,ind_B)=0.003564.
% За 0<q<1: Cov=0 <=> a_1=a_0 или b_1=b_0.
% P(E|A)=10/19, P(E|B)=25/52, P(E|A\cap B)=250/277.
\end{problemstar}

\begin{problemstar}[Урна с неизвестен състав]
За фиксирано $n\ge2$ избираме равномерно $K$ от $\{0,1,\dots,n\}$.
След това използваме урна с $K$ червени и $n-K$ сини топки.
От нея се правят две тегления. Разгледайте отделно моделите с връщане
и без връщане. Нека $R_i$ е събитието ,,при $i$-тото теглене е изтеглена
червена топка``, за $i\in\{1,2\}$.
\begin{enumeratea}
  \item Ако първата топка е червена, докажете, че в двата модела
  \[
  \Pbb(K=k\mid R_1)=\frac{2k}{n(n+1)},
  \qquad
  \text{за }k\in\{0,\dots,n\}.
  \]

  \item Намерете $\Pbb(R_2\mid R_1)$ в двата модела и докажете, че
  резултатите са съответно
  \[
  \frac{2n+1}{3n}
  \qquad\text{и}\qquad
  \frac23.
  \]

  \item Обяснете вероятностно защо резултатите са различни за крайно $n$
  и защо разликата изчезва при $n\to\infty$.
\end{enumeratea}
% С връщане: sum_{k=0}^n 2k^2/(n^2(n+1)).
% Без връщане: sum_{k=0}^n 2k(k-1)/(n(n+1)(n-1)).
% Разликата е 1/(3n).
\end{problemstar}

\begin{researchtask}[Кога повторният тест не дава независимо доказателство?]
Нека $D\sim\Ber(0.01)$ показва наличие на заболяване, а $B\sim\Ber(q)$ --
проблемна лабораторна партида, като $D$ и $B$ са независими.
Всички тестове на даден човек се извършват с една и съща партида.
Условно на $(D,B)$ резултатите $T_1,T_2,\dots$ са независими и еднакво
разпределени, като $T_i=1$ означава положителен резултат.

Първоначално нека $q=0.10$, а вероятностите за положителен резултат са
зададени чрез
\[
\begin{array}{c|cc}
 & B=0 & B=1\\ \hline
\Pbb(T_i=1\mid D=1,B) & 0.95 & 0.80\\
\Pbb(T_i=1\mid D=0,B) & 0.05 & 0.30
\end{array}
\]
\begin{enumeratea}
  \item Изведете точна формула за
  \[
  p_k=\Pbb(D=1\mid T_1=\cdots=T_k=1),
  \qquad
  \text{за }k\ge1.
  \]

  \item Сравнете я с наивния модел, който първо осреднява чувствителността
  и вероятността за фалшиво положителен резултат по $B$, а след това приема,
  че резултатите от тестовете са условно независими при дадено $D$.
  Обяснете защо това допускане може да промени отговора.

  \item Направете таблица за $k=1,\dots,5$. Проверете, че
  \[
  p_1=\frac{17}{152},
  \qquad
  p_2=\frac{701}{1592}\approx0.4403,
  \]
  докато наивният модел дава приблизително $0.6109$ при $k=2$.

  \item Разгледайте параметричната решетка
  \[
  q\in\{0,0.05,0.10,0.20\},
  \qquad
  \Pbb(T_i=1\mid D=0,B=1)\in\{0.10,0.20,0.30,0.40\},
  \]
  при запазени останали параметри. Намерете при кои двойки наивната оценка
  за $p_2$ е завишена с поне $0.10$.

  \item Винаги ли наивният модел завишава $p_2$, ако позволим да се променят
  и четирите вероятности в таблицата? Дайте доказателство или контрапример
  и обяснете полученото.
    \item Винаги ли наивният модел завишава $p_2$, ако позволим да се променят
  и четирите вероятности в таблицата? Намерете необходимо и достатъчно
  условие наивната оценка да е завишена, дайте контрапример и
  интерпретирайте условието.    
\end{enumeratea}
% За k=1,...,5, точен / наивен модел:
% 0.1118/0.1118, 0.4403/0.6109, 0.7472/0.9514,
% 0.9055/0.9959, 0.9680/0.9997.
%
% Наивна минус точна оценка за p_2:
% редове q; колони P(T=1|D=0,B=1).
% q=0:       0      0      0      0
% q=0.05: 0.007  0.058  0.130  0.199
% q=0.10: 0.013  0.090  0.171  0.223
% q=0.20: 0.022  0.110  0.161  0.171
%
% Разлика >=0.10 при:
% (0.05,0.30), (0.05,0.40), (0.10,0.30), (0.10,0.40),
% (0.20,0.20), (0.20,0.30), (0.20,0.40).
%
% Контрапример за последната точка:
% q=0.10; чувствителности 0.05,0.95;
% вероятности за фалшиво положителен резултат 0.05,0.05.
% Точна оценка: 37/136 ≈ 0.2721.
% Наивна оценка: 196/2671 ≈ 0.0734.
\end{researchtask}

\noindent\textbf{Lean.} Нека $A,B\subseteq\Omega$, а индикаторите приемат
стойности в $\mathbb N$. Използвайте дефиницията от първата тема и докажете
в Lean 4 с Mathlib, че за всяко $x\in\Omega$
\[
\ind_{A\cap B}(x)=\ind_A(x)\ind_B(x).
\]

% !TeX program = pdflatex
% Седмица 3. Фрагмент за включване в основния сборник.
% Използва неговата преамбюла и означения.

\section{Случайни величини и дискретни разпределения}

\subsection*{Основни задачи}

\begin{problembook}
Играч залага 5 лева и хвърля два независими стандартни зара. При две шестици получава 100 лева, а
при точно една шестица -- 5 лева. При нито една шестица не получава нищо.
Нека $X$ е чистата (нетната) печалба на играча. Намерете закона и (математическото)
очакването на $X$. Справедлива ли е играта? 
% E[изплащане] = 100/36 + 50/36 = 25/6; нетно -5/6 лв.; справедлив залог 25/6 ≈ 4.17 лв.
\end{problembook}

\begin{problembook}
Човек хвърля честна монета и при ези прави крачка напред, а при тура --
назад. Резултатите от хвърлянията са независими. Каква е вероятността след
10 хвърляния да се намира:
\begin{enumeratea}
  \item в началната точка;
  \item на разстояние 2 от нея;
  \item пет крачки пред нея?
\end{enumeratea}
% (a) C(10,5)/2^10 = 63/256; (b) 2C(10,4)/2^10 = 105/256; (c) 0 -- позицията е четна.
\end{problembook}

\begin{problembook}
Нека $Z_1$ и $Z_2$ са резултатите от хвърлянето на два независими зара и
$X=Z_1+Z_2$.
\begin{enumeratea}
  \item Заровете са стандартни. Намерете закона, очакването и дисперсията
  на $X$.
  \item Нека
  \[
  \Pbb(Z_i=1)=\Pbb(Z_i=6)=\frac14
  \quad\text{и}\quad
  \Pbb(Z_i=j)=\frac18,
  \qquad\text{за }i\in\{1,2\},\ j\in\{2,3,4,5\}.
  \]
  Намерете закона, очакването и дисперсията
  на $X$ и покажете два начина за намиране на очакването.
  \item За стандартни зарове нека $Y=|X-7|$. Намерете закона на $Y$ и
  $\Ebb Y$. Сравнете $\Ebb Y$ с $|\Ebb X-7|$ и $\Ebb Y^2$ с $\Var X$.
\end{enumeratea}
% (a) E=7, Var=35/6. (b) E=7 (линейност или от закона), Var=15/2.
% (c) P(Y=0)=6/36, P(Y=k)=2(6-k)/36; E Y=35/18 != 0 = |E X-7|: очакването
% не комутира с функция на X. E Y^2=35/6=Var X, защото Y^2=(X-7)^2 и E X=7.
\end{problembook}

\begin{problem}
В кутия има 7 лампи, от които 3 са дефектни. Избират се равномерно 4
различни лампи.
Намерете закона и очакването на броя избрани дефектни лампи. Намерете
очакването и чрез индикатори, без да сумирате целия закон. Намерете и
дисперсията чрез индикатори. Независими ли са използваните индикатори?
% Закон (1,12,18,4)/35; E=12/7; Var=24/49; индикаторите са отрицателно корелирани.
\end{problem}

\begin{problembook}
Два независими стандартни зара се хвърлят пет пъти, като резултатите от
различните хвърляния също са независими. Каква е вероятността
броят на хвърлянията със сума 6 да бъде точно 2? Намерете очакването и
дисперсията на този брой.
% Bin(5,5/36): C(5,2)(5/36)^2(31/36)^3 ≈ 0.1232; E=25/36; Var=775/1296.
\end{problembook}

\begin{problempencil}[Липса на памет; К1, СЕМ 2023]
Нека $X\sim\Geom(p)$, т.е. $\Pbb(X=k)=p(1-p)^k$ за $k\in\{0,1,2,\dots\}$, където
$p\in(0,1]$.
\begin{enumeratea}
  \item Докажете, че за всички цели числа $m,n\ge0$, за които $\Pbb(X\ge n)>0$,
  \[
  \Pbb(X\ge m+n\mid X\ge n)=\Pbb(X\ge m).
  \]
  \item Обратно, нека $Y$ е случайна величина със стойности в
  $\{0,1,2,\dots\}$, за която горното равенство е изпълнено за всички
  цели числа $m,n\ge0$ с $\Pbb(Y\ge n)>0$. Докажете, че $Y\sim\Geom(p)$ за някое
  $p\in(0,1]$.
\end{enumeratea}
% Нека h(n)=P(Y>=n). Тогава h(m+n)=h(m)h(n). Ако h(1)=0, Y=0 п.с.;
% иначе h(n)=q^n, където q=h(1)<1, и p=1-q.
\end{problempencil}

\begin{problembook}
Да приемем, че броят на слабите земетресения в даден район за $t$ месеца
има Поасоново разпределение с параметър $2t$. Каква е вероятността за три
месеца да има по-малко от четири земетресения?
% lambda=6: 61 e^{-6} ≈ 0.151.
\end{problembook}

\begin{probleminteresting}[St. Petersburg paradox]
Казино иска вход $A$ лева. Хвърляме честна монета до първото ези, като
резултатите от хвърлянията са независими. Ако то се появи при $n$-тото
хвърляне, наградата е $2^n$ лева. Намерете очакваната
награда. За кои $A$ лично бихте играли? Какво се променя, ако наградата е
$\min\{2^n,2^{30}\}$ лева? Как сравнявате получения
резултат с отговора си за стойностите на $A$, при които бихте играли?
% E=∞; при ограничена награда E=31 лв.
\end{probleminteresting}

\subsection*{Допълнителни задачи}

\begin{problem}
В урна има 5 бели и 3 черни топки. Теглим до появата на бяла. Нека $X$ е
броят на изтеглените черни топки. Намерете закона, очакването и дисперсията
на $X$ при:
\begin{enumeratea}
  \item теглене без връщане;
  \item теглене с връщане.
\end{enumeratea}
% (a) (35,15,5,1)/56; E=1/2; Var=15/28.
% (b) Geom(5/8) върху {0,1,...}: E=(1-p)/p=3/5; Var=(1-p)/p^2=24/25.
\end{problem}

\begin{problem}
$80\%$ от домашните принтери работят без допълнителни настройки, като
резултатите за различните принтери и седмици са независими. Фирма продава
10 принтера седмично. Намерете вероятността поне 9 да работят веднага;
вероятността това да се случи през пет поредни седмици; вероятността първата
седмица, в която по-малко от 9 принтера работят веднага, да е $k$-тата.
Пресметнете последната вероятност при $k=3$ и при $k=21$ и намерете
очаквания брой седмици до първата такава седмица.
% q=P(Bin(10,0.8)>=9)=0.37581; q^5=0.00750; T е геометрична върху {1,2,...}:
% P(T=k)=q^{k-1}(1-q); P(T=3)≈0.0882; P(T=21)≈1.97e-9; E T=1/(1-q)≈1.602.
\end{problem}

\begin{problempencil}[К1, СЕМ 2023]
\begin{enumeratea}
  \item Нека $X_1,\dots,X_n$ са независими и еднакво разпределени с
  $F(m)=\Pbb(X_1\le m)$, и нека $X_{\min}=\min\{X_1,\dots,X_n\}$ и
  $X_{\max}=\max\{X_1,\dots,X_n\}$. Изразете
  $\Pbb(X_{\min}\le m)$ и $\Pbb(X_{\max}\le m)$ чрез $F(m)$ и $n$.
  \item Хвърлят се 5 стандартни зара. Във всеки кръг хвърляме едновременно
  всички останали зарове и отстраняваме тези, показали шестица.
  Всички хвърляния са независими. Нека $K$ е броят на кръговете до
  отстраняването на всички зарове. Какво е разпределението на броя
  хвърляния на един отделен зар? Намерете $\Pbb(K\le n)$ и $\Pbb(K=n)$
  за всяко цяло число $n\ge1$.
\end{enumeratea}
% (a) 1-(1-F(m))^n; F(m)^n.
% (b) отделен зар: геометрично върху {1,2,...} с p=1/6;
% P(K<=n)=(1-(5/6)^n)^5; P(K=n)=(1-(5/6)^n)^5-(1-(5/6)^{n-1})^5.
\end{problempencil}

\begin{problem}
Нека $X$ е случайна величина със стойности в $\{0,1,2,\dots\}$ и крайно
очакване.
\begin{enumeratea}
  \item Докажете, че
  \[
  \Ebb X=\sum_{k=1}^\infty\Pbb(X\ge k)=\sum_{n=0}^\infty\Pbb(X>n).
  \]
  \item Проверете формулата за $X\sim\Geom(p)$.
  \item Намерете $\Ebb K$ за броя на кръговете $K$ от предходната задача и
  го сравнете с очаквания брой хвърляния на един отделен зар.
\end{enumeratea}
% (a) Разменете реда на сумиране в sum_k k P(X=k) = sum_k sum_{j=1}^k P(X=k).
% (b) P(X>=k)=(1-p)^k; sum_{k>=1}(1-p)^k=(1-p)/p.
% (c) E K = sum_{n>=0} (1-(1-(5/6)^n)^5) = sum_{j=1}^5 (-1)^{j+1} C(5,j)/(1-(5/6)^j)
%     = 3698650986/283994711 ≈ 13.02 кръга срещу 6 хвърляния за един зар.
\end{problem}

\begin{problem}
Провеждат се независими опити на Бернули с вероятност за успех $p$, където
$0<p<1$. Намерете вероятността $r$-тият успех да настъпи точно на
$(k+r)$-тия опит, където $r\ge1$ и $k\ge0$ са цели числа.
Нека $X_i$ е индикаторът за успех в $i$-тия опит. Запишете събитието чрез
$X_i$ и частичните суми $S_n=\sum_{i=1}^nX_i$, като $S_0=0$.
% C(k+r-1,r-1) p^r (1-p)^k; {S_{k+r-1}=r-1, X_{k+r}=1}.
\end{problem}

\begin{problem}
Вероятността един изстрел да улучи е $0.001$. Направени са 5000 независими
изстрела и са нужни поне две попадения. Намерете точната вероятност и
Поасоновото приближение. Колко голяма е числената разлика?
% Точно 0.95964; Poisson(5): 1-6e^{-5}=0.95957; разлика ≈ 7e-5.
\end{problem}

\begin{problem}
Подводница стреля $n$ пъти; изстрелите са независими и всяко торпедо улучва
с вероятност $p$. При попадение един от $m$-те отсека се избира равномерно и
независимо от останалите попадения. Намерете вероятността поне два различни
отсека да бъдат наводнени.
% 1 + (m-1)(1-p)^n - m(1-p+p/m)^n.
\end{problem}

\begin{problem}[Кибритените кутии на Банах]
Пушач носи две кутии с по $n$ клечки. Всеки път избира независимо една от
тях с вероятност $1/2$ и взема клечка. Когато за първи път посегне към празна
кутия, каква е вероятността в другата да са останали точно
$k\in\{0,\dots,n\}$ клечки?
% C(2n-k,n) 2^{-(2n-k)}, k=0,...,n.
% Отрицателно биномно разсъждение (вж. задачата за r-тия успех): последният
% избор е (n+1)-вият на празната кутия, преди него има n-k избора на другата;
% множителят 2 е за избора коя кутия се оказва празна.
\end{problem}

\subsection*{Предизвикателства и самостоятелно изследване}


\begin{problemstar}
Нека $\pi$ е равномерна случайна пермутация на $\{1,\dots,n\}$ и
\[
I_n=\#\{(i,j):1\le i<j\le n,\ \pi(i)>\pi(j)\}.
\]
Докажете, че
\[
\Ebb I_n=\frac{n(n-1)}4
\qquad\text{и}\qquad
\Var(I_n)=\frac{n(n-1)(2n+5)}{72}.
\]
% I_n=sum_{k=1}^n U_k, където U_k са независими и равномерни върху {0,...,k-1};
% E U_k=(k-1)/2, Var U_k=(k^2-1)/12.
% Защо U_k са независими и равномерни: равномерна пермутация на [n] се получава,
% като n се вмъкне на равномерно избрана позиция в равномерна пермутация на [n-1];
% новите инверсии са точно елементите вдясно от n -- равномерно върху {0,...,n-1}
% и независимо от пермутацията на [n-1]. Индукция по n.
% Директно с индикатори: за i<j<k ковариациите са +1/12,+1/12,-1/12.
\end{problemstar}

\begin{problemstar}[Coupon collector]
Нека $X_1,X_2,\dots$ са независими и еднакво разпределени, равномерни върху
$\{1,\dots,m\}$, където $m\ge2$, и
\[
T=\inf\{n\ge1:\{X_1,\dots,X_n\}=\{1,\dots,m\}\}.
\]
\begin{enumeratea}
  \item Намерете $\Pbb(T\le n)$ за $n\ge1$.
  \item Докажете, че
  \[
  \Ebb T=mH_m,
  \qquad
  \Var(T)=m^2H_m^{(2)}-mH_m,
  \]
  където $H_m=\sum_{j=1}^m\frac1j$ и
  $H_m^{(2)}=\sum_{j=1}^m\frac1{j^2}$.
\end{enumeratea}
Къде използвате независимостта на времената за чакане и къде е достатъчна
само линейността на очакването?
% (a) sum_{j=0}^m (-1)^j C(m,j)(1-j/m)^n.
% (b) T е сума от независими геометрични величини върху {1,2,...}
% с p_k=(m-k+1)/m, т.е. E=1/p_k, Var=(1-p_k)/p_k^2.
% Независимостта на тегленията дава геометричните закони и независимостта
% на времената; за сумата на очакванията стига линейност, а за дисперсията
% използваме независимостта. При m=6: E=14.7, Var=38.99.
\end{problemstar}

\begin{researchtask}[Обобщен събирач на купони]
Нека $X_1,X_2,\dots$ са независими и еднакво разпределени случайни
величини със стойности в $[m]=\{1,\dots,m\}$, където $m\ge2$, и
\[
\Pbb(X_n=i)=p_i>0,
\qquad\text{за }i\in[m],
\qquad\text{като }\sum_{i=1}^m p_i=1.
\]
Нека $T$ е броят на тегленията до появата на всички $m$ стойности.
\begin{enumeratea}
  \item Докажете, че
  \[
  \Ebb T=
  \sum_{\emptyset\ne A\subseteq[m]}
  \frac{(-1)^{|A|+1}}{\sum_{i\in A}p_i}
  \qquad\text{и}
\qquad
  \Ebb T=
  \int_0^\infty
  \left(1-\prod_{i=1}^m\bigl(1-e^{-p_it}\bigr)\right)\,dt.
  \]
  Обяснете защо първата формула е неподходяща за пряко числено пресмятане
  при $m=30$. Предложете числен метод за пресмятане на интеграла и
  начин да проверите точността му.
  \item Сравнете стойностите от интегралната формула със симулация при
  $m=10,30,100$ за следните модели:
  \begin{enumerate}[label=(\roman*)]
    \item равномерни вероятности:
    $p_i=1/m$ за $i\in[m]$;
    \item вероятности
    \[
    p_i=\frac{i^{-\alpha}}{\sum_{j=1}^m j^{-\alpha}},
    \qquad\text{за }i\in[m],
    \]
    като разгледате $\alpha=1$ и $\alpha=2$;
    \item един рядък купон:
    \[
    p_1=m^{-2}
    \quad\text{и}\quad
    p_i=\frac{1-m^{-2}}{m-1},
    \qquad\text{за }i\in\{2,\dots,m\}.
    \]
  \end{enumerate}
  Посочете броя на симулациите и обсъдете точността на полученото средно.
  Обяснете кой механизъм определя мащаба на $\Ebb T$ във всеки модел.
\end{enumeratea}
Можете да използвате AI за доказателството и числения метод, като
проверите резултатите и обясните как сте ги проверили.
% Интегралът следва и без поасонизация: включване--изключване за P(T>n),
% сумиране по n>=0 (формулата E T = sum_{n>=0} P(T>n) от допълнителните задачи),
% после 1/p(A)=integral_0^infty exp(-p(A)t)dt.
% Разкриването на крайното произведение дава интегралната формула.
% E T (интеграл), m=10/30/100: равномерно 29.3/119.8/518.7 (=mH_m);
% alpha=1: 56.2/318.3/1856.5 (асимптотично m (log m)^2);
% alpha=2: 250/3274/50870 (ред на растеж m^2 log m, много редки купони);
% един рядък: 103.5/907.7/10014 (асимптотично 1/p_1=m^2).
\end{researchtask}

\noindent\textbf{Lean.} Нека $S$ е крайно непразно множество, $X,Y:S\to\mathbb R$
и $a,b\in\mathbb R$. За крайното равномерно вероятностно пространство върху $S$
положете
\[
\Ebb X=\frac1{|S|}\sum_{s\in S}X(s).
\]
Формализирайте в Lean 4 с Mathlib и докажете линейността
\[
\Ebb(aX+bY)=a\,\Ebb X+b\,\Ebb Y.
\]
Докажете първо, по индукция по $S$ и без да използвате
\texttt{Finset.sum\_add\_distrib}, че
\[
\sum_{s\in S}\bigl(X(s)+Y(s)\bigr)
=
\sum_{s\in S}X(s)+\sum_{s\in S}Y(s);
\]
за останалото са полезни \texttt{Finset.mul\_sum} и \texttt{Finset.sum\_div}.
Защо непразността на $S$ е нужна за дефиницията на $\Ebb$, но не и за
равенството за сумите?

\section{Зависимост, суми и пораждащи функции}

\subsection*{Част I. Съвместни разпределения и условно очакване}

\begin{problembook}
Хвърляме независимо два червени и един син стандартен зар. Нека $X$ е броят
шестици върху червените зарове, а $Y$ --- броят двойки върху трите зара.
Намерете съвместния закон и $\Pbb(X>0\mid Y=1)$. Независими ли са $X$ и $Y$?
Можете ли да отговорите на последния въпрос и без да пресмятате целия
съвместен закон?
\end{problembook}

\begin{problembook}
Честна монета се хвърля независимо четири пъти. Нека $X$ е броят езита в
първите три хвърляния, а $Y$ --- броят езита в последните две.
\begin{enumeratea}
  \item Намерете съвместния закон, условните закони, $\Pbb(X=Y)$ и
  $\Pbb(X>1\mid Y=1)$.
  \item Намерете $\Ebb(X\mid Y)$ и $\Ebb(Y\mid X)$ и техните закони.
  Проверете, че $\Ebb[\Ebb(X\mid Y)]=\Ebb X$.
  \item Намерете $\Cov(X,Y)$ по два начина: чрез съвместния закон и чрез
  представяне на $X$ и $Y$ като суми от индикатори.
\end{enumeratea}
\end{problembook}

\begin{problembook}
Нека $X$ е равномерно разпределена върху $\{-1,0,1\}$ и $Y=X^2$.
\begin{enumeratea}
  \item Покажете, че $\Cov(X,Y)=0$, но $X$ и $Y$ са зависими.
  \item Нека сега $X$ и $Y$ приемат стойности в $\{0,1\}$.
  Докажете, че $\Cov(X,Y)=0$ е еквивалентно на независимостта им.
  Защо това не противоречи на предишния пример?
  \item За случайни величини $X$ и $Y$ с крайни положителни дисперсии
  докажете, че $|\Cor(X,Y)|\le1$. Кога се достига равенство?
\end{enumeratea}
\end{problembook}

\begin{problempencil}
В една кутия има 4 бели и 6 черни топки, а в друга --- 7 бели и 3 черни.
От всяка кутия изваждаме равномерно по една топка, независимо от другата.
Към двете извадени топки добавяме една бяла и поставяме трите в трета кутия.
\begin{enumeratea}
  \item Каква е вероятността равномерно избрана топка от третата кутия
  да е бяла?
  \item Теглим от третата кутия с връщане до първата бяла топка. При фиксиран
  състав на кутията тегленията са независими и равномерни.
  Нека $X$ е броят тегления. Намерете $\Ebb X$.
  Има ли $X$ геометрично разпределение?
  \item Намерете $\Pbb(X>k)$ за всяко цяло число $k\ge0$ и докажете, че
  $\Pbb(X>k+1\mid X>k)$ расте с $k$.
  Какво научаваме за състава на кутията след поредица от черни топки?
\end{enumeratea}
\end{problempencil}

\subsection*{Допълнителни}

\begin{problembook}
От $\{1,2,3,4,5\}$ избираме равномерно три числа без връщане. Нека $X$ е
средното по големина, а $Y$ --- най-малкото.
\begin{enumeratea}
  \item Намерете съвместния и маргиналните закони.
  Независими ли са $X$ и $Y$?
  \item Намерете $\Cov(X,Y)$ и $\Cor(X,Y)$.
  \item Намерете закона, очакването и дисперсията на $X-2Y$.
\end{enumeratea}
\end{problembook}

\begin{problempencil}
Нека $(X_1,\ldots,X_n)$ е равномерна случайна пермутация на
$\{1,\ldots,n\}$, където $n\ge2$, и $S=X_1+\cdots+X_n$.
\begin{enumeratea}
  \item Намерете $\Ebb S$, $\Var S$, $\Ebb X_i$ и $\Var X_i$.
  \item Докажете формулата за дисперсия на сума и я използвайте, за да
  намерите $\Cov(X_i,X_j)$ за $i\ne j$. Обяснете знака на ковариацията.
  \item За $1\le k\le n$ намерете $\Var(X_1+\cdots+X_k)$.
  Сравнете с дисперсията при $k$ независими тегления с връщане.
\end{enumeratea}
\end{problempencil}

\begin{problempencil}
На всяка страна на празен зар независимо записваме равномерно число от 1 до 6.
След това хвърляме получения зар $n\ge2$ пъти. Падналите страни са независими,
равномерно избрани и независими от надписите. Нека $T_k$ е числото, което
виждаме при $k$-тото хвърляне.
\begin{enumeratea}
  \item Какъв е законът на $T_1$? Намерете $\Cor(T_1,T_2)$.
  Независими ли са $T_1$ и $T_2$?
  \item Ако $X=T_1+\cdots+T_{n-1}$ и $Y=T_2+\cdots+T_n$,
  намерете $\Cor(X,Y)$.
  \item Какво се променя в двете корелации, ако надписите са независими и
  равномерни върху $\{199,\ldots,999\}$? А ако надписите са независими
  и имат един и същ произволен закон с крайна положителна дисперсия?
\end{enumeratea}
\end{problempencil}

\begin{problem}[Урна на P\'olya]
В урна има една червена и една черна топка. На всяка стъпка избираме равномерно
една топка, връщаме я и добавяме още една от същия цвят. Нека $X_n$ е
индикаторът за червена топка при $n$-тото теглене и $W_n=X_1+\cdots+X_n$.
Намерете $\Pbb(X_n=1)$, $\Cov(X_i,X_j)$ за $i\ne j$ и закона на $W_4$.
Какъв закон предполагате за $W_n$?
\end{problem}

\subsection*{Предизвикателства}

\begin{problemstar}[Независимост по двойки]
Нека $q$ е просто число, а $A$ и $B$ са независими и равномерно разпределени
върху $\{0,\ldots,q-1\}$. Всички следващи сметки са по модул $q$. За
$t\in\{0,\ldots,q-1\}$ полагаме
\[
X_t=A+tB\pmod q,\qquad X_\infty=B.
\]
\begin{enumeratea}
  \item Докажете, че величините $X_0,\ldots,X_{q-1},X_\infty$ са равномерно
  разпределени и независими по двойки, но не са съвместно независими.
  Какво определят стойностите на две различни координати?
  \item Нека $N$ е броят нулеви координати в това семейство.
  Намерете закона, очакването и дисперсията на $N$.
  \item Пресметнете дисперсията и чрез индикатори.
  Достатъчна ли е независимостта по двойки, за да се събират дисперсиите?
\end{enumeratea}
\end{problemstar}

\begin{problemstar}
За случайни величини $X$ и $Y$ с крайни втори моменти и случайна величина $Z$
докажете, че
\[
\Cov(X,Y)=\Ebb\!\left[\Cov(X,Y\mid Z)\right]
+\Cov\!\left(\Ebb[X\mid Z],\Ebb[Y\mid Z]\right).
\]
Каква формула получавате при $Y=X$?
Нека сега $Z\sim\Ber(1/2)$, $U$ е равномерно разпределена върху $\{-1,1\}$,
$Z$ и $U$ са независими, и
\[
X=3Z+U,\qquad Y=3Z-U.
\]
\begin{enumeratea}
  \item Намерете условните ковариация и корелация при фиксирано $Z$.
  \item Намерете $\Cov(X,Y)$ и $\Cor(X,Y)$.
  \item Как обяснявате разликата в знаците? Свържете отговора си с двата
  члена във формулата.
\end{enumeratea}
\end{problemstar}

\subsection*{Част II. Пораждащи функции и случайни суми}

\begin{problempencil}[К1, КН 2024]
За неотрицателна целочислена случайна величина $X$ означаваме
$g_X(s)=\Ebb[s^X]$, $0\le s\le1$.
\begin{enumeratea}
  \item Нека $X$ е равномерно разпределена върху $\{1,\ldots,6\}$.
  Намерете $g_X$ и изразете $\Ebb X$ и $\Var X$ чрез нейни производни.
  \item Ако $X$ и $Y$ са независими неотрицателни целочислени случайни
  величини, докажете, че $g_{X+Y}=g_Xg_Y$.
  \item[$(\mathrm{c})^\star$] Вярно ли е обратното твърдение?
\end{enumeratea}
\end{problempencil}

\begin{problembook}
Нека $X$ и $Y$ са независими. Чрез пораждащи функции намерете закона на $X+Y$,
ако $X\sim\Bin(n,p)$ и $Y\sim\Bin(m,p)$, а след това --- ако
$X\sim\Poi(\lambda)$ и $Y\sim\Poi(\mu)$.
Намерете по същия начин закона на $G_1+\cdots+G_r$, където
$G_1,\ldots,G_r$ са независими и $G_i-1\sim\Geom(p)$, т.е. всяка $G_i$
брои опитите до първия успех.
Защо равенството на пораждащите функции определя закона еднозначно?
Може ли сумата на независими $\Bin(n,p)$ и $\Bin(m,q)$ величини да е
биномично разпределена, ако $n,m\ge1$, $0<p,q<1$ и $p\ne q$?
\end{problembook}

\begin{problempencil}
Броят посетители на стадион е $N\sim\Poi(\lambda)$. Всеки посетител избира
равномерно един от десетте входа, независимо от останалите и от $N$.
\begin{enumeratea}
  \item Намерете закона, очакването и дисперсията на броя посетители през
  първия вход.
  \item Нека по-общо $N,X_1,X_2,\ldots$ са независими неотрицателни
  целочислени случайни величини, като $(X_i)$ имат един и същ закон.
  За $S=X_1+\cdots+X_N$, със $S=0$ при $N=0$, докажете, че
  \[
  g_S(s)=g_N\bigl(g_{X_1}(s)\bigr).
  \]
  Ако $N$ и $X_1$ имат крайни втори моменти, намерете $\Ebb S$ и $\Var S$.
  \item Приложете полученото към задачата за стадиона.
  Получете закона от (а) и чрез формулата за пълната вероятност.
\end{enumeratea}
\end{problempencil}

\begin{problembook}
$A$ и $B$ играят независими партии. Във всяка партия $A$ печели с вероятност
$2/3$, а $B$ --- с вероятност $1/3$. Играта спира, когато някой спечели две
последователни партии. Нека $T$ е броят изиграни партии.
Намерете $\Ebb T$ и вероятността играта да завърши с победа на $A$.
Каква информация за миналото е достатъчна, за да опишем продължението на играта?
Намерете и пораждащата функция на $T$.
\end{problembook}

\subsection*{Допълнителни}

\begin{problempencil}
Два отбора изпълняват по една дузпа във всеки кръг. Всички изпълнения са
независими, с вероятност за гол $0.75$ за първия отбор и $0.80$ за втория.
Изпълняват се всичките пет редовни кръга. При равенство след тях играта
продължава, докато в един кръг единият отбор отбележи, а другият не.
\begin{enumeratea}
  \item Каква е вероятността във всеки от първите пет кръга двата отбора
  да имат еднакъв резултат от изпълнението си?
  \item Каква е вероятността играта да продължи след петия кръг?
  А след десетия?
  \item Намерете очаквания общ брой изпълнени дузпи.
\end{enumeratea}
\end{problempencil}

\begin{problem}
Лотарийните билети са номерирани от 000000 до 999999.
Намерете вероятността равномерно избран билет:
\begin{enumeratea}
  \item да има сума на цифрите 21;
  \item да има равни суми на първите и последните три цифри;
  \item да има сума на първите три цифри с 2 по-голяма от сумата на последните три.
\end{enumeratea}
Достатъчно е да изразите отговорите чрез коефициенти на полином.
\end{problem}

\begin{problempencil}
На $N\sim\Poi(\lambda)$ души се обаждаме с предложение. Човек $i$ би се
съгласил при $K_i$-тото обаждане, където $K_1,K_2,\ldots$ са независими,
еднакво разпределени и независими от $N$, с
$\Pbb(K_i=n)=f(n)$ за $n\ge1$ и $\sum_{n\ge1}f(n)=1$.
Нека $X_n$ е броят хора, които биха се съгласили при $n$-тото обаждане.
\begin{enumeratea}
  \item Докажете, че величините $(X_n)_{n\ge1}$ са независими и
  $X_n\sim\Poi(\lambda f(n))$.
  \item Спираме при първия кръг без ново съгласие. Ако
  $T=\inf\{n\ge1:X_n=0\}$ и $S=X_1+\cdots+X_T$, докажете, че
  \[
  \Ebb S=\lambda\Ebb[F(T)],\qquad F(k)=\sum_{j=1}^k f(j).
  \]
  Защо тук не можем направо да приложим формулата за очакване на случайна сума
  с брой събираеми, независим от самите събираеми?
\end{enumeratea}
\end{problempencil}

\begin{problempencil}
Два стандартни зара се хвърлят независимо. Можем ли да изпишем други
положителни цели числа върху страните им, така че при независими хвърляния
законът на сумата да остане същият? Повторения на надписите са позволени.
Намерете конструкция, различна от двата стандартни зара.
\end{problempencil}

\subsection*{Предизвикателства}

\begin{problemstar}\label{prob:week4-penney}
Хвърляме независимо честна монета и записваме резултатите с $H$ и $T$.
За дума $w=w_1w_2w_3$ нека $\tau_w$ е номерът на хвърлянето, при което тя
се появява за първи път като три последователни резултата.
Разгледайте думите $a=HHT$ и $b=THH$.
\begin{enumeratea}
  \item Изберете състояния, които пазят достатъчно информация за появата на
  двете думи, и начертайте възможните преходи.
  \item Намерете $\Pbb(\tau_a<\tau_b)$ и $\Ebb(\tau_a\wedge\tau_b)$.
  \item Докажете, че $\tau_a\wedge\tau_b<\infty$ почти сигурно,
  без да използвате пресметнатото очакване.
\end{enumeratea}
\end{problemstar}

\begin{problemstar}[Разклоняващ процес]
Една популация започва от един индивид. Всеки индивид има 0 деца с вероятност
$1-p$ и 2 деца с вероятност $p$, където $0<p<1$, независимо от всички останали
индивиди във всички поколения. Нека $Z_n$ е броят индивиди в поколение $n$,
$Z_0=1$, и
\[
f(s)=1-p+ps^2,\qquad T=\inf\{n\ge0:Z_n=0\}.
\]
\begin{enumeratea}
  \item Докажете, че $g_{Z_n}=f^{\circ n}$.
  \item Покажете, че $q_n=\Pbb(T\le n)=f^{\circ n}(0)$ и че $(q_n)$ расте.
  \item Докажете, че $q=\Pbb(T<\infty)$ е най-малкото решение на $f(q)=q$
  в $[0,1]$, и го намерете в зависимост от $p$.
  \item Намерете $\Ebb Z_n$. При $p=1/2$ то е постоянно, а популацията
  изчезва почти сигурно. Как обяснявате това?
\end{enumeratea}
\end{problemstar}

\subsection*{Самостоятелно изследване}

\begin{researchtask}[Играта на Penney]
Двама играчи избират по една различна дума от $\{H,T\}^3$.
След това хвърляме независимо честна монета, докато се появи една от двете думи.
Печели играчът, чиято дума се е появила първа.
\begin{enumeratea}
  \item Обобщете модела от задача~\ref{prob:week4-penney} и напишете програма,
  която за всяка от 56-те наредени двойки $(a,b)$ пресмята точно
  \[
  p(a,b)=\Pbb(\tau_a<\tau_b),\qquad
  e(a,b)=\Ebb(\tau_a\wedge\tau_b).
  \]
  Използвайте рационална аритметика и проверете ръчно поне една двойка.
  \item За всяка първа дума $a$ намерете най-добрите отговори $b$ на втория
  играч. Има ли предимство играчът, който избира втори?
  Формулирайте общо правило за най-добрия отговор и го докажете.
  \item Начертайте ориентиран граф с ребро $a\to b$, когато $b$ е най-добър
  отговор на $a$. Намерете цикъл и обяснете какво показва той за сравнението
  между думите.
\end{enumeratea}
\end{researchtask}

\noindent\textbf{Lean.} Формализирайте в Lean 4 с Mathlib и докажете, че
за крайни множества $S,T$ и функции $a:S\to\mathbb R$ и $b:T\to\mathbb R$,
\[
\sum_{s\in S}\sum_{t\in T}a(s)b(t)
=
\left(\sum_{s\in S}a(s)\right)
\left(\sum_{t\in T}b(t)\right).
\]
Обяснете как това тъждество се използва при доказателството на
$g_{X+Y}=g_Xg_Y$ за независими неотрицателни целочислени случайни
величини $X$ и $Y$, приемащи краен брой стойности.
\end{document}